
\documentclass[11pt,a4paper]{article}

\usepackage[utf8]{inputenc}
\usepackage[T1]{fontenc}
\usepackage{lmodern}
\usepackage{microtype}
\usepackage{amsmath,amssymb,amsthm}
\usepackage{mathtools}
\usepackage{geometry}
\usepackage{hyperref}

\geometry{a4paper,margin=25mm}

\newtheorem{theorem}{Theorem}[section]

\title{\textbf{Identity Bedrock: A Minimal Cryptographic Substrate and Its Systemic Implications}}
\author{Jean-François Bouchard}
\date{March 2026}

\begin{document}

\maketitle

\begin{abstract}
We present Identity Bedrock, a deployed identity system constructed from a minimal cryptographic substrate grounded in one-way functions (OWFs).

The system separates identity cryptography from directory consensus via a chain-agnostic append-only binding layer.

The system has been deployed in production for over five years, with publicly verifiable records.

We show that naming, content resolution, execution, connectivity, and economic interaction emerge from this identity substrate.
\end{abstract}

\section{Introduction}

This work documents a deployed system.

\section{Architecture}

\subsection{Identity}

Users control identity via keypairs $(sk, pk)$.

\subsection{Directory}

An append-only substrate binds identifiers to public keys.

\section{Security}

\begin{theorem}
Authentication is unforgeable assuming signature security and directory integrity.
\end{theorem}

\section{Derived Properties}

\subsection{Naming}

\begin{equation}
(n, v, pk, \sigma), \quad \sigma = \text{Sign}_{sk}(n \parallel v)
\end{equation}

\subsection{Execution}

\begin{equation}
Execution = Hash(Input, Environment, Command)
\end{equation}

\section{Deployment}

\url{https://explorer.emercoin.com/dpo/PrivateNESS.Network}

\section{Implementation}

EmerDNS, NVS, AuxPoW, Fiber.

\section{Conclusion}

Minimal identity substrate is sufficient.

\section*{Acknowledgments}

AI-assisted editing.

\begin{thebibliography}{3}

\bibitem{rompel1990} Rompel.
\bibitem{rfc1035} Mockapetris.
\bibitem{nakamoto2008} Nakamoto.

\end{thebibliography}

\end{document}
